%%% ===============================================================================
\chapter{First Real Chapter addressing first Research Problem}\label{sec:problem1}
%%% ===============================================================================

\emph{Purpose and scope}.
After you stated research context (SM1), research problem (SM2), and research method (SM3) in \zcref{sec:introduction} and possibly \zcref{sec:problem-exposition}, the remainder of your entire report addresses execution (SM4), results (SM5), and interpretation (SM6).
Do this by addressing the sub-problems through scientific arguments that follow the 6 steps SM1-SM6.

Have a short chapter introduction that recalls and explains the first research problem of your thesis.
The problem has to show up in the introduction in \zcref{sec:introduction:research-question} or in \zcref{sec:problem-exposition:research-problems} already.
This provides the background (SM1) for this chapter while the first research problem of the thesis becomes the research question/hypothesis (SM2) for this chapter.

Next, explain in the chapter intro how you solve the research problem in this chapter by breaking it down in further sub-problems.
By this, you outline the method (SM3) through which you are going to solve the problem of this chapter.
This is necessary to give the reader guidance of what's to come in this chapter and how it fits into the thesis as a whole.
Explain that you will address the first sub-problem in \zcref{sec:problem1:subproblem1} and the second sub-problem in \zcref{sec:problem1:subproblem2}, etc. The sections then provide the details for execution and results.

% ---------------------------------------------------------------------------------
\section{First Sub-Problem}\label{sec:problem1:subproblem1}
% ---------------------------------------------------------------------------------

\emph{The first paragraph describes the first sub-problem and develops the requirements a solution has to satisfy (SM2 for this section).}
The requirements have to be based on the knowledge and reasoning developing in the preceding chapters and sections.
Try to use an example to illustrate the problem and the desired properties of the solution.
Check that every term/concept you use here has already been defined already in a previous section.
If you cannot describe your problem without defining new terms, you may have to add another section before this one that develops the terms and concepts you need to explain the problem.

\emph{The second paragraph describes the method/approach how you address the problem (SM3 for this section).}
Describe the method in a level of detail that allows another student to reproduce your steps.
Make use of appendices % (see \zcref{sec:appendix1})
if certain details take too much space.

\emph{The third, fourth, and following paragraph provides details on applying the method or developing a new approach, i.e., execution (SM4) and may explain results (SM5)}, i.e., details on the steps needed to reproduce the results.

Results (SM5) come in different forms, e.g., conceptual diagrams, algorithms, tables, charts, a list of articles from a literature research etc. You must reference them (e.g., ``\zcref{fig:my_label} shows the runtime'') and describe the results in text.
If you use diagrams, tables, or charts, you cannot expect the reader to know what to you expect them to see in a diagram, table or chart.
Describe to them how to read these, explain the meaning of particular elements, point out special observations.
Describe the results, but do not interpret them.
Make use of appendices if certain details take too much space.

\begin{figure}
  \centering
  %%%\includegraphics{/path/to/figure.pdf}
  \caption{A scientific figure that has to be explained in the text}
  \label{fig:my_label}
\end{figure}

\emph{After describing the results, you may interpret them (SM6).}
Here you can infer what a particular observation means (for you), how it can be applied, or what others can do with it.
Do not interpret your results before you have described all of them.
This is a common mistake done by most beginner writers.
Get to the point: the final finding or interpretation.
But you forget that your reader does not understand yet what you are interpreting - they do not know yet what you do know.
Interpret the results only after describing all of them.
The interpretation must be based on the written description only.
Then you can be sure that your readers can follow your interpretation and reach the same conclusions as you have.

Ideally, your interpretation leads to the next sub-problem in \zcref{sec:problem1:subproblem2}.

% ---------------------------------------------------------------------------------
\section{Second Sub-Problem}\label{sec:problem1:subproblem2}
% ---------------------------------------------------------------------------------

You now build on the solution to the first sub-problem of \zcref{sec:problem1:subproblem1} (SM1) and recall second sub-problem (SM2, you detailed in the introduction of this chapter) and follow the same pattern as before (SM3-SM6).

Note that not all sections may not include all parts SM1-SM6 in all detail.
Some sections need not restate the background (SM1) or the research problem (SM2) if a previous section already defined them.
Sometimes, a section is only dedicated to describing the method (SM3) and execution (SM4) and does not contain any results or interpretations.
Sometimes results (SM5) and interpretations (SM6) only come in the evaluation chapter.

What is important for you when you are writing a scientific argument is not to slavishly have SM1-SM6 in each section explicitly, but that you are always fully aware of the following:
%
\begin{itemize}
  \item Which step of a scientific argument am I currently writing (SM1, SM2, ..., SM6)?
  \item Does the step that I am writing come in the right order, i.e., if you are writing about execution (SM4, e.g., details of building a model), is there a preceding paragraph or section that describes the method (SM3) and is that one preceded by a clear statement of the (sub-)problem addressed (SM2)?
  \item Are you really \emph{not} writing interpretation SM6 before SM5, SM4, or SM3?
  \item Is it clear to the reader which part of the scientific argument you are currently making?
\end{itemize}
